% GENERATED FILE. Edit canonical lessons, then run npm run generate:books.
% canonical-source-hash: fnv1a64:585cad64dcfb5ffa
% canonical-lessons: LA-C382-senator, LA-C382-patricius, LA-C382-plebeius, LA-C382-nobilis, LA-C382-popularis

\chapter{Senators and People}
\label{ch:la-senators-and-people}
\hlchaptermodality{382}

\begin{chapteropening}
\textbf{By the end of this chapter:} \emph{I can talk about senators, patricians and plebeians, the nobility and the popular side.}

Complete the last lesson of chapter 382: I can talk about senators, patricians and plebeians, the nobility and the popular side.
\end{chapteropening}

\section[senātor, senātōris]{senātor, senātōris — a senator}

\label{lesson:LA-C382-senator}



\begin{quote}
Before the new one: say the Latin for a copyist, then the Latin for a land surveyor.
\end{quote}

\subsection*{You'll want to know: senātor, senātōris}
\textbf{senātor, senātōris} — \textquotedblleft{}a senator\textquotedblright{}.

\begin{quote}
Senātōrēs in cūriā sedent.
\end{quote}

\begin{quote}
\emph{The senators sit in the senate house.}
\end{quote}

\emph{senātor, senātōris} m., a member of the \emph{senātus}, the council of elders (\emph{senex}, an old man).

\begin{cousinweb}
\emph{senātor} is built on \emph{senātus} \textquotedblleft{}the council of elders\textquotedblright{} (EtymDB), which comes from \emph{senex} \textquotedblleft{}an old man\textquotedblright{} (Lewis \& Short). English \emph{senator} comes from it.
\end{cousinweb}

\begin{grammarlens}[title={What you've built}]
One more word for telling, explaining and describing at B1.
\end{grammarlens}

\subsection*{Your turn}
\begin{itemize}
\raggedright
  \item \emph{Say it:} \emph{senātor, senātōris}
  \item \emph{Say it:} \emph{senātor, senātōris}, once more
  \item \emph{Say it:} a sentence of your own with \emph{senātor, senātōris}
  \item \emph{Recall:} say the Latin for a copyist, then the Latin for a land surveyor, then say \emph{senātor, senātōris} again
\end{itemize}

\begin{tcolorbox}[breakable,colback=teal!4,colframe=teal!35!black,title={Before you move on}]
What does senātor, senātōris mean? (\textquotedblleft{}A senator\textquotedblright{}.) Say it once more.
\end{tcolorbox}
\section[patricius, patricia, patricium]{patricius, patricia, patricium — patrician, of the old noble families}

\label{lesson:LA-C382-patricius}



\begin{quote}
Before the new one: say the Latin for a land surveyor, then the Latin for a senator.
\end{quote}

\subsection*{You'll want to know: patricius, patricia, patricium}
\textbf{patricius, patricia, patricium} — \textquotedblleft{}patrician, of the old noble families\textquotedblright{}.

\begin{quote}
Familia eius patricia et antīqua est.
\end{quote}

\begin{quote}
\emph{His family is patrician and ancient.}
\end{quote}

\emph{patricius, patricia, patricium}, built on \emph{patrēs}, the fathers, an old name of the senators. English \emph{patrician}.

\begin{cousinweb}
\emph{patricius} is built on \emph{patrēs} \textquotedblleft{}the fathers\textquotedblright{} (EtymDB; Lewis \& Short), the plural of \emph{pater} used of the first senators. English \emph{patrician} comes from it.
\end{cousinweb}

\begin{grammarlens}[title={What you've built}]
One more word for telling, explaining and describing at B1.
\end{grammarlens}

\subsection*{Your turn}
\begin{itemize}
\raggedright
  \item \emph{Say it:} \emph{patricius, patricia, patricium}
  \item \emph{Say it:} \emph{patricius, patricia, patricium}, once more
  \item \emph{Say it:} a sentence of your own with \emph{patricius, patricia, patricium}
  \item \emph{Recall:} say the Latin for a land surveyor, then the Latin for a senator, then say \emph{patricius, patricia, patricium} again
\end{itemize}

\begin{tcolorbox}[breakable,colback=teal!4,colframe=teal!35!black,title={Before you move on}]
What does patricius, patricia, patricium mean? (\textquotedblleft{}Patrician, of the old noble families\textquotedblright{}.) Say it once more.
\end{tcolorbox}
\section[plēbēius, plēbēia, plēbēium]{plēbēius, plēbēia, plēbēium — plebeian, of the common people}

\label{lesson:LA-C382-plebeius}



\begin{quote}
Before the new one: say the Latin for a senator, then the Latin for patrician.
\end{quote}

\subsection*{You'll want to know: plēbēius, plēbēia, plēbēium}
\textbf{plēbēius, plēbēia, plēbēium} — \textquotedblleft{}plebeian, of the common people\textquotedblright{}.

\begin{quote}
Prīmus cōnsul plēbēius ā populō laudātus est.
\end{quote}

\begin{quote}
\emph{The first plebeian consul was praised by the people.}
\end{quote}

\emph{plēbēius, plēbēia, plēbēium}, built on \emph{plēbs}, the common people. English \emph{plebeian}.

\begin{cousinweb}
\emph{plēbēius} is built on \emph{plēbs, plēbis} \textquotedblleft{}the common people\textquotedblright{} (EtymDB; Lewis \& Short). English \emph{plebeian} comes from it.
\end{cousinweb}

\begin{grammarlens}[title={What you've built}]
One more word for telling, explaining and describing at B1.
\end{grammarlens}

\subsection*{Your turn}
\begin{itemize}
\raggedright
  \item \emph{Say it:} \emph{plēbēius, plēbēia, plēbēium}
  \item \emph{Say it:} \emph{plēbēius, plēbēia, plēbēium}, once more
  \item \emph{Say it:} a sentence of your own with \emph{plēbēius, plēbēia, plēbēium}
  \item \emph{Recall:} say the Latin for a senator, then the Latin for patrician, then say \emph{plēbēius, plēbēia, plēbēium} again
\end{itemize}

\begin{tcolorbox}[breakable,colback=teal!4,colframe=teal!35!black,title={Before you move on}]
What does plēbēius, plēbēia, plēbēium mean? (\textquotedblleft{}Plebeian, of the common people\textquotedblright{}.) Say it once more.
\end{tcolorbox}
\section[nōbilis, nōbile]{nōbilis, nōbile — well known; noble}

\label{lesson:LA-C382-nobilis}



\begin{quote}
Before the new one: say the Latin for patrician, then the Latin for plebeian.
\end{quote}

\subsection*{You'll want to know: nōbilis, nōbile}
\textbf{nōbilis, nōbile} — \textquotedblleft{}well known; noble\textquotedblright{}.

\begin{quote}
Nōbilis vir est, sed pauper.
\end{quote}

\begin{quote}
\emph{He is a man of noble birth, but poor.}
\end{quote}

\emph{nōbilis, nōbile}: first known (it is built on \emph{nōscō}, to get to know), then of a known family. The \emph{nōbilēs} had a consul among their ancestors. English \emph{noble}.

\begin{cousinweb}
\emph{nōbilis} is built on \emph{nōscō}, old \emph{gnōscō} \textquotedblleft{}to get to know\textquotedblright{} (EtymDB; Lewis \& Short), with the ending \emph{-bilis}: one who can be known. English \emph{noble} comes from it through French.
\end{cousinweb}

\begin{grammarlens}[title={What you've built}]
One more word for telling, explaining and describing at B1.
\end{grammarlens}

\subsection*{Your turn}
\begin{itemize}
\raggedright
  \item \emph{Say it:} \emph{nōbilis, nōbile}
  \item \emph{Say it:} \emph{nōbilis, nōbile}, once more
  \item \emph{Say it:} a sentence of your own with \emph{nōbilis, nōbile}
  \item \emph{Recall:} say the Latin for patrician, then the Latin for plebeian, then say \emph{nōbilis, nōbile} again
\end{itemize}

\begin{tcolorbox}[breakable,colback=teal!4,colframe=teal!35!black,title={Before you move on}]
What does nōbilis, nōbile mean? (\textquotedblleft{}Well known; noble\textquotedblright{}.) Say it once more.
\end{tcolorbox}
\section[populāris, populāre]{populāris, populāre — of the people, popular; a fellow countryman}

\label{lesson:LA-C382-popularis}



\begin{quote}
Before the new one: say the Latin for plebeian, then the Latin for well known.
\end{quote}

\subsection*{You'll want to know: populāris, populāre}
\textbf{populāris, populāre} — \textquotedblleft{}of the people, popular; a fellow countryman\textquotedblright{}.

\begin{quote}
Ōrātor populāris plēbem ad sē trahēbat.
\end{quote}

\begin{quote}
\emph{The popular speaker drew the common people to himself.}
\end{quote}

\emph{populāris, populāre}. The \emph{populārēs} were politicians who worked through the people's assemblies. As a noun, \emph{populāris} is a fellow countryman. English \emph{popular}.

\begin{cousinweb}
\emph{populāris} is built on \emph{populus} \textquotedblleft{}the people\textquotedblright{} (EtymDB; Lewis \& Short), with the ending \emph{-āris} \textquotedblleft{}belonging to\textquotedblright{}. English \emph{popular} comes from it.
\end{cousinweb}

\begin{grammarlens}[title={What you've built}]
One more word for telling, explaining and describing at B1.
\end{grammarlens}

\subsection*{Your turn}
\begin{itemize}
\raggedright
  \item \emph{Say it:} \emph{populāris, populāre}
  \item \emph{Say it:} \emph{populāris, populāre}, once more
  \item \emph{Say it:} a sentence of your own with \emph{populāris, populāre}
  \item \emph{Recall:} say the Latin for plebeian, then the Latin for well known, then say \emph{populāris, populāre} again
\end{itemize}

\begin{tcolorbox}[breakable,colback=teal!4,colframe=teal!35!black,title={Before you move on}]
What does populāris, populāre mean? (\textquotedblleft{}Of the people, popular; a fellow countryman\textquotedblright{}.) Say it once more.
\end{tcolorbox}
